\documentclass[12pt,a4paper]{article}

% ---------- Pacotes ----------
\usepackage[utf8]{inputenc}
\usepackage[T1]{fontenc}
\usepackage[brazilian]{babel}
\usepackage{amsmath,amssymb}
\usepackage{graphicx}
\usepackage{xcolor}
\usepackage{enumitem}
\usepackage{fancyhdr}
\usepackage{lastpage}
\usepackage{titlesec}
\usepackage{listings}
\usepackage{tikz}
\usepackage{pgfplots}
\usepackage[hidelinks]{hyperref}
\usepackage[top=3.2cm,bottom=3cm,left=2.5cm,right=2.5cm,headheight=40pt]{geometry}

\pgfplotsset{compat=1.16}

% ---------- Cores ----------
\definecolor{unbblue}{RGB}{0,0,255}      % azul dos títulos
\definecolor{unbdark}{RGB}{20,70,110}    % azul escuro do cabeçalho
\definecolor{unbgreen}{RGB}{0,140,70}

% ---------- Títulos das seções (azul, sem numeração) ----------
\titleformat{\section}{\sffamily\bfseries\large\color{unbblue}}{}{0pt}{}
\titlespacing*{\section}{0pt}{1.8em}{0.8em}

% ---------- Listas com (a), (b), ... em azul ----------
\setlist[enumerate,1]{label=\textcolor{unbblue}{(\alph*)},leftmargin=2em,itemsep=0.8em}
\setlist[itemize]{label=\textemdash,itemsep=0.1em}

% ---------- Código ----------
\lstset{
  basicstyle=\ttfamily\small,
  breaklines=true,
  columns=fullflexible,
  keepspaces=true,
  literate={á}{{\'a}}1 {ã}{{\~a}}1 {ç}{{\c c}}1 {ê}{{\^e}}1 {é}{{\'e}}1 {õ}{{\~o}}1 {í}{{\'i}}1 {ú}{{\'u}}1 {ô}{{\^o}}1
}

% ---------- Cabeçalho e rodapé ----------
\pagestyle{fancy}
\fancyhf{}
\renewcommand{\headrulewidth}{0.4pt}
\renewcommand{\footrulewidth}{0.4pt}

\lhead{%
  \raisebox{-0.2\height}{%
  \begin{tikzpicture}[scale=0.11]
    % logotipo estilizado (substitua por \includegraphics{logo_unb.png} se tiver o arquivo)
    \fill[unbdark] (0,0) -- (4,0) -- (4,3) -- cycle;
    \fill[unbgreen] (0,0) -- (0,3) -- (4,3) -- cycle;
    \fill[unbdark] (5,0) -- (9,0) -- (5,3) -- cycle;
    \fill[unbgreen] (9,0) -- (9,3) -- (5,3) -- cycle;
  \end{tikzpicture}}%
  \hspace{0.3em}{\sffamily\bfseries\Large\color{unbdark}Universidade de Brasília}}
\rhead{%
  {\sffamily\large FCTE}\;%
  \raisebox{-0.25\height}{\tikz\node[circle,fill=unbgreen,text=white,inner sep=2pt]{\Large$\gamma$};}}

\lfoot{\small\sffamily Faculdade de Ciências e Tecnologias em Engenharia -- FCTE\\
       \textcolor{unbblue}{\href{http://www.fcte.unb.br}{www.fcte.unb.br}}}
\rfoot{\small\sffamily\thepage/\pageref{LastPage}}

\setlength{\parindent}{1.5em}

\begin{document}

% ---------- Título ----------
\begin{center}
  {\sffamily\bfseries\large\color{unbblue} Processamento de Sinais -- 02/2026}\\[2.2em]
  {\sffamily\bfseries\large\color{unbblue} Lista de Exercícios 1}\\[0.8em]
  {\sffamily\bfseries\large\color{unbblue} Conceitos básicos}
\end{center}

\vspace{0.5em}
\noindent\textbf{Nome:} Henrique de Oliveira Bernardes \hfill \textbf{Data:} \today\\
\textbf{Matrícula:} 231011471

\section{Tópicos tratados}
\begin{itemize}
  \item Conceito de sinais em tempo discreto e em tempo contínuo.
  \item Relação entre sinais em tempo discreto e em tempo contínuo.
  \item Amostragem a taxas constantes.
  \item Teorema de Nyquist\\
        (transformada de Fourier do produto de um sinal por um trem de impulsos).
  \item Reconstrução de Shannon a partir de amostras de sinais.
  \item Quantização de sinais.
  \item Diferenças entre processamento analógico e processamento digital de sinais.
\end{itemize}

% =====================================================
\section{Problema 1}
Acerca de conceitos básicos de processamento de sinais, responda as perguntas que se seguem.

\begin{enumerate}
  \item O que é um sinal?

  \item O que é um sinal em domínio contínuo e um sinal em domínio discreto?

  \item Suponha que um sinal em tempo contínuo $x_c$ é amostrado a um intervalo regular. Qual o
        nome dado a este intervalo e ao inverso dele, no contexto de amostragem de sinais?

  \item Com respeito ao item anterior, qual a equação que relaciona o sinal em tempo contínuo $x_c$
        com o sinal em tempo discreto $x$?

  \item Com respeito ao item c: há alguns critérios que, quando respeitados, garantem que é possível
        calcular, teoricamente sem erro, o valor do sinal em tempo contínuo em qualquer instante, a
        partir das amostras adquiridas a intervalos regulares. Qual desses critérios é o mais usado?
        Trata-se de um critério necessário, suficiente, ou os dois?

  \item Supondo que o critério do item acima é respeitado, há uma equação em forma fechada para
        cálculo de um valor do sinal em tempo contínuo, num instante arbitrário $t$, a partir das
        amostras? Ou seja, há uma fórmula para calcular $x_c(t)$ a partir dos vários $x[n]$'s? Se há,
        forneça essa fórmula.
\end{enumerate}

% =====================================================
\section{Problema 2}
Defina um sinal arbitrário em tempo contínuo, pela sobreposição de 3 senoides de frequências
$f_1 = 200$ hertz, $f_2 = 400$ hertz, $f_3 = 600$ hertz e amplitudes $A_1 = 1$, $A_2 = 2$ e
$A_3 = 1{,}5$, respectivamente. Use fases diferentes para cada senoide. A respeito desse sinal,
responda os seguintes itens.

\begin{enumerate}
  \item Elabore um programa em MATLAB ou Octave para geração do gráfico de $x_c$, em um intervalo
        de tempo correspondente a 4 períodos do sinal. No gráfico, indique claramente o
        significado dos eixos horizontal e vertical. Forneça o código e o gráfico obtido. Use uma
        taxa de amostragem alta o suficiente para que não se observe o aspecto discreto implícito
        no gráfico.

  \item Escreva um programa em MATLAB ou Octave para amostrar o sinal $x_c$ a uma taxa de 3 kHz.
        Trace o gráfico das 30 primeiras amostras do sinal. Utilize representação de cada amostra
        por uma barra vertical e círculo na extremidade da barra (gráfico do tipo \textit{stem}).

  \item A respeito do item anterior, determine o intervalo de tempo contínuo, em segundos,
        correspondente às 30 amostras representadas (considere o intervalo entre a primeira e a
        última amostra).
\end{enumerate}

% =====================================================
\section{Problema 3}
Considere que um sinal $x_c:\mathbb{R}\to\mathbb{R}$ tem a transformada de Fourier ilustrada no
gráfico abaixo (os valores da transformada são nulos fora do intervalo de frequências mostrado no
gráfico).

\begin{center}
\begin{tikzpicture}
\begin{axis}[
  width=0.8\textwidth, height=7cm,
  xmin=-2000, xmax=2000, ymin=0, ymax=12500,
  xtick={-2000,-1500,-1000,-500,0,500,1000,1500,2000},
  ytick={0,2000,4000,6000,8000,10000,12000},
  xlabel={Frequência $f$ (hertz)},
  ylabel={$\hat{x}_c(f)$}, ylabel style={rotate=-90},
  grid=both, grid style={gray!40},
  scaled ticks=false
]
\addplot[blue!70!black, thick] coordinates {(-2000,0) (-1000,0) (0,12000) (1000,0) (2000,0)};
\end{axis}
\end{tikzpicture}
\end{center}

Com base nessa informação, responda os seguintes itens.

\begin{enumerate}
  \item Forneça a equação para a transformada de Fourier do sinal dado por
        \[
          y_c(t) = \sum_{n=-\infty}^{\infty} x_c(t)\,\delta(t - nT_s),
        \]
        onde $\delta$ é a função impulso unitário (ou delta de Dirac). Considere que
        $T_s = 0{,}4$ ms.

  \item Trace o gráfico da transformada de Fourier do sinal $y_c$ definido no item anterior.

  \item Repita o item \textbf{a)}, considerando agora $T_s = 0{,}625$ ms.

  \item Repita o item \textbf{b)}, considerando agora $T_s = 0{,}625$ ms.

  \item Explique a importância do teorema que embasa, na teoria de processamento de sinais, a
        resposta aos itens anteriores.
\end{enumerate}

% =====================================================
\section{Problema 4}
Suponha que um sinal em tempo contínuo $x_c$ foi amostrado a um intervalo constante de 0,5 ms.
Suponha ainda que essa condição de amostragem está de acordo com o critério de Nyquist, para o
sinal em questão.

O gráfico abaixo mostra o sinal em tempo discreto obtido, que é nulo fora da faixa de valores de
tempo mostrada.

\begin{center}
\begin{tikzpicture}
\begin{axis}[
  width=0.7\textwidth, height=7cm,
  xmin=-3.5, xmax=5.5, ymin=-3, ymax=8,
  xtick={-3,-2,-1,0,1,2,3,4,5},
  ytick={-3,-2,-1,0,1,2,3,4,5,6,7,8},
  xlabel={$n$}, ylabel={$x[n]$}, ylabel style={rotate=-90},
  grid=both, grid style={gray!40}
]
\addplot[ycomb, blue!70!black, mark=*, thick] coordinates
  {(-3,0) (-2,2) (-1,3) (0,1) (1,-1) (2,1) (3,2) (4,7) (5,1)};
\draw[red] (axis cs:-3,0) -- (axis cs:5,0);
\end{axis}
\end{tikzpicture}
\end{center}

Com base nessas informações:

\begin{enumerate}
  \item Calcule o valor do sinal em tempo contínuo no instante 0.
  \item Calcule o valor do sinal em tempo contínuo no instante 0,5 ms.
  \item Calcule o valor do sinal em tempo contínuo no instante 0,75 ms.
  \item Calcule o valor do sinal em tempo contínuo no instante 1 ms.
  \item Calcule o valor do sinal em tempo contínuo no instante 1,25 ms.
  \item Com auxílio do MATLAB, Python ou outro programa/linguagem, trace um gráfico do sinal em
        tempo contínuo entre os instantes de $-2$ milissegundos e de 2,5 milissegundos.
\end{enumerate}

% =====================================================
\section{Problema 5}
O arquivo \texttt{amostras\_exemplo.mat} fornece 20 amostras de um sinal amostrado a uma taxa
constante de 4 kHz. A respeito desse sinal, responda os seguintes itens. Os índices de tempo são
representados pelo vetor $n_1$, enquanto que os valores do sinal naqueles índices são
representados pelo vetor $x_1$.

Considere que as amostras obtidas à taxa de 4 kHz são nulas, com exceção das 20 amostras dadas.
Considere ainda que a primeira amostra da lista de 20 corresponde ao índice de tempo $n = 0$.

\begin{enumerate}
  \item Determine o intervalo de tempo representado pelas 20 amostras.

  \item Supondo que a amostragem atende ao critério de Nyquist, trace um gráfico do sinal em tempo
        contínuo, em uma janela de tempo que inclua as 20 amostras e permita a visualização de
        todas elas, além do gráfico em tempo contínuo. Utilize o MATLAB, Python, ou outro
        programa/linguagem. No gráfico, apresente as amostras individuais (representadas como
        \textit{stems}) sobrepostas ao sinal em tempo contínuo. Note que é necessário calcular os
        instantes em s ou ms para cada amostra (e não representá-las, nesse caso, em índices
        inteiros de tempo).

  \item Se interpolarmos o sinal com outra técnica que não a de Shannon (por exemplo, usando
        interpolação polinomial ou por \textit{splines}), obteremos um sinal distinto, mas que
        também interpola as amostras. Em que sentido o sinal obtido pela equação de Shannon traduz
        mais fielmente o sinal real? Por quê?

  \item Com respeito ao item anterior, que característica matemática diferencia o sinal obtido
        pela equação de Shannon de qualquer outro sinal que interpole as amostras?
\end{enumerate}

% =====================================================
\section{Problema 6}
Quais são as duas etapas no processo de digitalização de um sinal em tempo contínuo? Descreva cada
etapa, e diga, em cada caso, se é possível em geral realizá-la sem perda teórica de informação ou
introdução de erro.

% =====================================================
\section{Problema 7}
Tanto a etapa de amostragem quanto a etapa de quantização podem afetar negativamente a informação
presente em um sinal originalmente analógico, quando feitas de forma não criteriosa.

No caso da amostragem, há teorias avançadas acerca de condições que, quando satisfeitas, garantem
que não há perda de informação no processo em questão.

No caso da quantização, o erro introduzido durante o processo pode ter implicações mais severas ou
menos severas para uma aplicação específica, dependendo do número de \textit{bits} utilizado.

Neste problema, aborda-se o efeito prático da alteração do número de níveis de quantização,
suposta linear, para o entendimento de um sinal de voz. Neste sentido, responda os seguintes itens.

\begin{enumerate}
  \item Utilizando o MATLAB ou o Python, grave um segmento de sua própria voz, com aproximadamente
        20 segundos de duração, ou adquira um exemplo de sinal de voz a partir de um áudio ou
        vídeo disponível publicamente (com o MATLAB, Python ou o Audacity, é possível extrair o
        sinal de áudio a partir de um arquivo de vídeo). A ideia é que o sinal represente um
        discurso falado, e não um canto ou manifestação musical. Registre bem a informação da taxa
        de aquisição e do número de bits utilizado na aquisição.

        No caso do MATLAB, é possível gravar um áudio (com um canal) de 20 segundos a 44,1 kHz e
        com 24 bits, utilizando o código a seguir.

\begin{lstlisting}
%  parametros de aquisicao
fs = 44100; % Taxa de amostragem em hertz
nBits = 24; % Número de bits por amostra
duracao = 20; % Duração em segundos

% Objeto para a gravação (o 1 se refere a um canal único de aquisição)
recObj = audiorecorder(fs, nBits, 1);
disp('Comece a fala.')

% Aquisição do áudio
recordblocking(recObj, duracao);
disp('Fim da gravação.');

% Gravação em arquivo
x = getaudiodata(recObj);
audiowrite('exemplo_voz.wav', x, fs)
\end{lstlisting}

  \item Ouça o sinal gravado. No MATLAB, você pode usar os comandos abaixo. Comente sobre a
        qualidade do áudio e sobre a inteligibilidade do que foi dito, a partir do sinal.

\begin{lstlisting}
[x, fs] = audioread('exemplo_voz.wav');
sound(x, fs)
\end{lstlisting}

  \item Escreva uma função que reduz o número de níveis de quantização de um sinal. A função
        recebe na entrada o sinal \texttt{x} e o número de bits \texttt{Nbits} correspondente à
        quantidade de níveis de quantização desejada. Em seguida, a função retorna um novo sinal,
        contendo no máximo $2^{\text{Nbits}}$ níveis distintos. Como dica, siga os passos indicados
        abaixo, em sua função.
        \begin{itemize}
          \item Normalize o sinal de entrada, colocando-o na faixa entre 0 e 1. Por exemplo:
\begin{lstlisting}
x = x - min(x);
x = x / max(x);
\end{lstlisting}
          \item Multiplique o sinal resultante por $2^{\text{Nbits}} - 1$. Os novos valores estarão
                entre 0 e $2^{\text{Nbits}} - 1$.
          \item Aproxime cada amostra para sua parte inteira (no MATLAB, isso é feito usando o
                comando \texttt{floor}). Assim, todos os valores serão inteiros entre 0 e
                $2^{\text{Nbits}} - 1$.
          \item Para ouvir o sinal obtido, conforme as orientações abaixo, use o comando
                \texttt{soundsc(x\_quantizado, fs)}.
        \end{itemize}

  \item Utilizando sua função, reduza o número de níveis de quantização para o correspondente a
        16 bits. Ouça o novo sinal. Comente sobre a qualidade e inteligibilidade.

  \item Utilizando sua função, reduza o número de níveis de quantização para o correspondente a
        8 bits. Ouça o novo sinal. Comente sobre a qualidade e inteligibilidade.

  \item Utilizando sua função, reduza o número de níveis de quantização para o correspondente a
        4 bits. Ouça o novo sinal. Comente sobre a qualidade e inteligibilidade.

  \item Utilizando sua função, reduza o número de níveis de quantização para o correspondente a
        2 bits. Ouça o novo sinal. Comente sobre a qualidade e inteligibilidade.

  \item Utilizando sua função, reduza o número de níveis de quantização para o correspondente a
        1 bit. Ouça o novo sinal. Comente sobre a qualidade e inteligibilidade.
\end{enumerate}

% =====================================================
\section{Problema 8}
Com que objetivos se realiza a digitalização de sinais analógicos? Levando em conta que a
amostragem de sinais pode levar a perda de informação, não seria melhor efetuar o processamento de
um sinal em sua forma analógica, em vez de digitalizá-lo?

% =====================================================
\section{Problema 9}
Cite e explique as vantagens e desvantagens do processamento digital de sinais em relação ao
processamento analógico de sinais.

% =====================================================
\section*{\sffamily\bfseries\color{unbblue} Referências}
\begin{enumerate}[label={[\arabic*]},leftmargin=2.5em,itemsep=0.6em]
  \item P. Diniz, E. A. B. da Silva, and S. L. Netto. \textit{Processamento Digital de Sinais:
        Projeto e Análise de Sistemas}. Bookman, Porto Alegre, RS, 2ª edição, 2014.
  \item R. G. Lyons. \textit{Understanding digital signal processing}. Prentice Hall, Upper Saddle
        River, NJ, 3\textsuperscript{rd} edition, 2010.
  \item S. K. Mitra. \textit{Digital Signal Processing: a Computer-Based Approach}. McGraw-Hill
        Education, 4\textsuperscript{th} edition, 2010.
  \item A. V. Oppenheim, R. W. Schafer, et al. \textit{Discrete-time signal processing}.
        Prentice-Hall signal processing series. Prentice Hall, Englewood Cliffs, N.J., 2\textsuperscript{nd}
        edition, 1999.
\end{enumerate}

\end{document}